% ─────────────────────────────────────────────────────────────────────────────
\section{Discussion}
\label{sec:discussion}
% ─────────────────────────────────────────────────────────────────────────────

\subsection{Interpretation and Implications}
\label{subsec:implications}

The results carry a clear operational message: the two halves of the diagnosis
problem are \emph{not} equally hard, and they should not be treated as a single
modeling task. For \textbf{Anomaly Detection}, a plain Logistic Regression on
184 global statistics reaches F1\,=\,0.691 with full interpretability and
millisecond inference. This is already a deployable safety screen: it can flag
roughly seven of every ten anomalous flights for human review using nothing more
than per-sensor means, spreads, and trends. The single most discriminative
channel---oil pressure (\texttt{E1 OilP})---matches the physical intuition that
lubrication state is the primary early indicator of engine health, which gives
maintenance engineers a defensible reason to trust the flag.

For \textbf{Fault Classification}, the picture is different and arguably more
important to state honestly. Even with feature selection and an Oracle
hierarchical grouping, the classical ceiling (F1\,=\,0.420) leaves the majority
of the 19 fault types poorly recognized, and several minority classes are never
recovered at all (Table~\ref{tab:perclass}). The end-to-end consequence is
concrete: of the flights our cascade gets wrong, \emph{40.6\%} are correctly
flagged as anomalous but assigned the \emph{wrong} fault type
(Table~\ref{tab:error_decomp}). Such an error does not endanger the aircraft---it
is still grounded for inspection---but it misdirects the maintenance action and
erodes trust in the system. The practical implication is that a global-statistics
classifier is fit for \emph{triage} (sick vs.\ healthy) but not yet for
\emph{prescription} (which part to fix).

Why does this matter beyond a class assignment? General aviation still runs on
calendar-based maintenance. A reliable AD screen alone---cheap, interpretable,
and trainable on a laptop---is enough to move part of the fleet toward
\emph{condition-based} inspection, catching emerging faults between scheduled
windows. The fault-typing problem, by contrast, is where temporal deep learning
earns its cost.

\subsection{Relation to Prior Work}
\label{subsec:related}

Our study sits downstream of three lines of NGAFID research. The dataset itself
was introduced by Yang and Desell \cite{ngafid2021} as a large-scale, real
(non-simulated) maintenance benchmark; we adopt their per-second 23-sensor
representation and physically isolated cross-validation protocol. Yang,
LaBella, and Desell \cite{convtok2023} then showed that convolutional
transformers operating on the raw sequence outperform shallow baselines for
predictive maintenance---consistent with our finding that the residual FC gap is
temporal in nature. Finally, the LMSD framework \cite{lmsd2025} proposed the very
two-stage decomposition (global screening for AD, micro-scale kernels for FC)
that we reproduce here with classical models.

An important caveat governs this comparison. The deep learning numbers we cite
(ConvTokMHSA F1\,=\,0.799 on AD; MMK~Net F1\,=\,0.623 on FC) come from the LMSD
preprint, which was \emph{withdrawn by its own authors} over acknowledged
``methodological flaws \ldots\ in the experimental validation and
metric-computation procedures.'' We deliberately retained the comparison because
the \emph{qualitative} story---global context helps AD, local micro-patterns help
FC---is corroborated independently by \cite{convtok2023} and by our own ablations
(global-attention ConvTokMHSA scoring \emph{below} flat XGBoost on FC). But the
precise gap magnitudes (9\,pp on AD, 20\,pp on FC) should be read as provisional.
If anything, the withdrawal strengthens the case for our classical baselines:
when a state-of-the-art result is retracted, a transparent, reproducible,
laptop-trainable system that reports honest cross-validated variance becomes more
valuable, not less.

\subsection{Survey of Tools Investigated}
\label{subsec:tools}

This project was deliberately built on the open-source Python scientific stack
emphasized in the course. Table~\ref{tab:tools} surveys the principal tools, the
role each played, and the practical strengths and limitations we observed.

\begin{table}[H]
    \centering
    \caption{Survey of tools investigated, with observed strengths and limitations.}
    \label{tab:tools}
    \begin{tabular}{p{2.4cm} p{4.1cm} p{6.0cm}}
        \toprule
        \textbf{Tool} & \textbf{Role in this project} & \textbf{Observed strength / limitation} \\
        \midrule
        NumPy & Vectorized statistics for the 2{,}048$\times$23 flight matrices; closed-form OLS slope. & {\itshape +} $\sim5\times$ faster slope than \texttt{polyfit}. {\itshape --} manual NaN handling. \\
        \addlinespace
        pandas & Loading, joining maintenance labels, fold bookkeeping. & {\itshape +} expressive munging. {\itshape --} memory-heavy at 11{,}446$\times$184. \\
        \addlinespace
        scikit-learn \cite{sklearn2011} & LR, LinearSVM, RF, KNN; pipelines, scaling, CV, metrics. & {\itshape +} uniform API, fold-safe preprocessing. {\itshape --} no native time-series shape features. \\
        \addlinespace
        XGBoost \cite{xgboost2016} & Best flat classical model on FC; gain-based feature selection (224$\to$60). & {\itshape +} handles nonlinear interactions and imbalance. {\itshape --} sensitive to long-tail minority classes. \\
        \addlinespace
        MiniRocket \cite{minirocket2021} & Repr.~G: random convolutional kernels for temporal shape. & {\itshape +} fast, near-deterministic. {\itshape --} $\sim$10k features, low interpretability, modest FC gain. \\
        \addlinespace
        tsfresh \cite{tsfresh2018} & Repr.~H: autocorrelation, FFT, entropy, nonlinearity descriptors. & {\itshape +} rich feature library. {\itshape --} expensive to compute, many redundant features. \\
        \addlinespace
        Matplotlib / Seaborn & EDA, class-distribution and PCA-variance figures. & {\itshape +} publication-quality output. {\itshape --} verbose for faceted plots. \\
        \addlinespace
        Keras / TensorFlow & Surveyed as the route to the temporal DL models that define the upper bound. & {\itshape +} the natural tool to close the FC gap. {\itshape --} out of scope here; left to future work. \\
        \bottomrule
    \end{tabular}
\end{table}

The central lesson from this tool survey is one of \emph{matching the
representation to the task}. The scikit-learn / XGBoost path is outstanding for
AD: fast, interpretable, and statistically honest under cross-validation. The
dedicated time-series transforms (MiniRocket, tsfresh) promised to recover the
temporal shape that FC needs, but in practice they inflated the feature space by
one to two orders of magnitude while yielding only marginal FC improvement---a
strong empirical signal that fault signatures are \emph{local} events that random
or generic features dilute. That negative result is what points unambiguously
toward locally restricted deep learning (Keras/TensorFlow), the one family of
tools we surveyed but did not train, as the principled next step.

\subsection{Limitations}
\label{subsec:limitations}

Three limitations bound our conclusions. First, all results use the 2-day
benchmark subset (11{,}446 flights, 19 classes); the full 36-class NGAFID is more
imbalanced and would likely lower every score. Second, the Oracle K\,=\,4 result
is an \emph{upper bound}: it assumes the coarse fault group is known at test time,
so the F1\,=\,0.420 figure is a ceiling, not a deployable number. Third, and most
importantly, our deep learning reference points inherit the reliability problem of
the withdrawn source preprint \cite{lmsd2025}; the comparison is directionally
sound but not numerically definitive.

\subsection{Future Work}
\label{subsec:future}

The error decomposition makes the priority unambiguous: \textbf{fault
classification is the bottleneck}, with twice the end-to-end leverage of anomaly
detection. The clearest next steps are (i) training a locally restricted temporal
model (MMK-style micro-kernel CNN in Keras/TensorFlow) on the raw sequences to
test whether the FC gap is genuinely temporal; (ii) re-running the full benchmark
once corrected LMSD results are published, to re-anchor the comparison on
trustworthy numbers; and (iii) addressing the 38:1 class imbalance directly
through resampling or cost-sensitive losses, since the long tail---not model
capacity alone---drives the minority-class failures observed in
Table~\ref{tab:perclass}.
